\documentclass[11pt]{article}
\usepackage[T1]{fontenc}
\usepackage{lmodern}
\usepackage[margin=1in]{geometry}
\usepackage{amsmath,amssymb,amsthm}
\usepackage{graphicx}
\usepackage[hidelinks]{hyperref}
\newtheorem*{theorem}{Published answer}
\setlength{\emergencystretch}{3em}
\title{Literature resolution of the general-tree\\$2$-Steiner distance-matrix rank problem}
\author{Run \texttt{fa\_banach\_001}, lane 18 (\texttt{agent\_lane\_18})}
\date{11 August 2026}

\begin{document}
\maketitle

\paragraph{Status.}
The open problem at the end of arXiv:2108.12798 was answered in full in 2022
by Ali Azimi and Sivaramakrishnan Sivasubramanian.  This packet records the
exact correspondence and claims no new mathematical result.

\paragraph{Definition.}
Let $T$ be a tree with vertex set $V(T)$.  If $S\subseteq V(T)$, its Steiner
distance $d_T(S)$ is the number of edges in the smallest connected subtree
containing $S$.  The $2$-Steiner distance matrix $D_2(T)$ is indexed by the
$2$-element subsets of $V(T)$, with
\[
 D_2(T)_{P,Q}=d_T(P\cup Q).
\]

\paragraph{Exact source question.}
Azimi, Bapat, and Goel prove in arXiv:2108.12798 that if $C$ is a caterpillar
on $N$ vertices with $p$ pendant vertices, then
\[
 \operatorname{rank}D_2(C)=2N-p-1.
\]
They end the paper with the following question.

\begin{center}
\includegraphics[width=.95\textwidth]{figures/open_problem_crop.png}
\end{center}

\paragraph{Exact later answer.}
The first theorem of Azimi and Sivasubramanian's 2022 paper treats an arbitrary
tree, not merely a caterpillar.

\begin{theorem}
Let $T$ be a tree on $n$ vertices having $p$ pendant vertices (leaves).  Then
\[
 \operatorname{rank}D_2(T)=2n-p-1.
\]
\end{theorem}

\begin{center}
\includegraphics[width=.88\textwidth]{figures/answer_theorem_crop.png}
\end{center}

Thus the caterpillar formula survives unchanged for every tree and answers
the source's literal request.

\paragraph{Proof route in the answer paper.}
The published proof is by induction on the number of vertices.  After the
three-vertex base case, it separates two reductions.  If a vertex supports at
least two leaves, deleting one leaf removes one row-space direction; if every
leaf has a degree-two neighbour, deleting a terminal two-vertex path removes
two directions.  The leaf count changes in exactly the corresponding way, so
the induction preserves the expression $2n-p-1$.  The proof also constructs
an explicit basis for the row space.  Later sections compute an inverse for
the nonsingular restriction to that basis, its determinant, and the inertia
of $D_2(T)$; these strengthen, but are not needed for, the rank answer.

\paragraph{Verification.}
The source PDF included here contains the displayed open problem on its page
23.  The answer PDF is the published version in \emph{Linear Algebra and its
Applications}; its abstract states the formula and Theorem~1 on journal page
67 gives it for every tree.  Its proof appears in Section~3.  The objects and
parameters agree exactly: $n$ is the order of the tree, $p$ is its number of
pendant vertices, and the matrix is indexed by vertex pairs with entries equal
to the Steiner distance of their union.

\paragraph{Conclusion and limitation.}
The target is classified as \texttt{literature\_already\_answered}.  No
novelty is claimed for the formula or proof.  Human review should verify only
the bibliographic attribution and the exact notational match recorded above.

\paragraph{References.}
\begin{enumerate}
\item A. Azimi, R. B. Bapat, and S. Goel, \emph{Steiner distance matrix of
caterpillar graphs}, Special Matrices 10 (2022), 167--188;
arXiv:2108.12798, doi:10.1515/spma-2022-0162.
\item A. Azimi and S. Sivasubramanian, \emph{The 2-Steiner distance matrix of
a tree}, Linear Algebra Appl. 655 (2022), 65--86,
doi:10.1016/j.laa.2022.09.007.
\end{enumerate}

\end{document}
